%==============================================================================
\section{Conclusion}
\label{sec:conclusion}
%==============================================================================

This thesis investigated whether methods independent of researcher-provided corpora can detect and mitigate hallucinations in LLM-generated Causal Loop Diagrams (CLDs), motivated by the reliability requirements of using LLMs for structured scientific knowledge generation. Across three research questions, we evaluated LLM-as-a-Judge configurations, LLM-as-a-correctors, uncertainty quantification (UQ) signals, automated literature-oriented validation via Deep Research, and adaptive deployment of Deep Research based on judge uncertainty signals.

Overall, the results support a nuanced conclusion: corpus-independent detection and correction methods can work in controlled settings, but their effectiveness degrades when the target is expert-grounded causal structure. In our dataset of three health-domain CLDs, correctness-judge performance is high in a controlled synthetic hallucination setting (GT Synth), but it drops in the literature ground-truth setting (GT Lit), with only the causal mechanistic prompt beating random baselines. In expert domains, disagreement with an expert CLD should not automatically be treated as fabrication: it can also reflect scope choices, tacit knowledge, or construct mismatch between an edge claim and what is explicitly evidenced in accessible sources. As a consequence, LLM-generated CLDs are best used as candidate drafts that focus expert attention on the contested edge space, rather than as self-validating outputs.

For corpus-independent hallucination signals, generator-side UQ metrics (perplexity, minimum probability, maximum window entropy, and cosine similarity) do not effectively flag hallucinations independently. They can contribute as features to supervised ensemble classifiers in-distribution, but they fail to generalize out-of-CLD, reinforcing their role as auxiliary signals for monitoring and triage rather than as standalone detectors.

For judge-guided correction, LLM-as-a-corrector consistently improves judge scores and slightly improves CLD quality in the synthetic setting, but not versus literature ground truth. This highlights a Goodhart-like risk in judge--corrector feedback loops: optimizing for proxy judge satisfaction without improving expert-grounded alignment.

Finally, for literature-oriented validation, we piloted a multi-agent Deep Research (DR) system as adaptive test-time compute. DR retrieves direct causal evidence for 62.4\% of false positive edges and 42.9\% of false negative edges. Human validation (\(n=50\)) suggests that at applicability threshold \(t=0.7\), 66\% of DR-evidenced edges pass the threshold and, among passing edges, 90.9\% are correct-causal. Under this calibration, extrapolated ground-truth enrichment yields substantial gains in aggregate edge F1: \(0.267 \to 0.527\) (Scenario 1: FP\(\to\)TP) and \(0.267 \to 0.582\) (Scenario 2: LLM+DR), with conservative estimates of 0.472 and 0.500. Using a mechanistic correctness judge to flag uncertain edges for adaptive DR deployment improves the cost--accuracy trade-off, delivering a 15\% accuracy-per-dollar increase.

%------------------------------------------------------------------------------
\subsection{Key Contributions}
\label{sec:conclusion_contributions}
%------------------------------------------------------------------------------

\begin{itemize}[noitemsep]
    \item We provide an end-to-end experimental evaluation of corpus-independent CLD generation and post-hoc validation, spanning generation, judging, correction, uncertainty signals, and evidence-oriented verification, with methods established in literature but not tested systematically in CLD context.
    \item We show that the \textit{Correctness Judge} performs strongly in controlled synthetic hallucination settings (GT Synth) but does not reliably transfer to literature ground truth (GT Lit), where only the Mechanistic prompt exceeds random baselines.
    \item We find that the \textit{Citation Judge} does not reliably discriminate valid from hallucinated edges under either GT Synth or GT Lit: retrieval and evidence asymmetries (notably, \textit{absence of evidence} being treated as \textit{evidence of absence} for negative edges) and systematic leniency in human validation make citation judging better suited to permissive screening than automated filtering.
    \item We benchmark generator-side uncertainty quantification (UQ) metrics as corpus-independent hallucination signals and show that they are weak as standalone detectors and do not generalize reliably across CLDs, but can serve as auxiliary features for monitoring and triage.
    \item We identify and characterize failure modes that matter specifically for CLDs (e.g., scope-vs-factual disagreement, plausibility bias, evidence-of-absence effects, and construct/applicability mismatch).
    \item We demonstrate that judge-guided correction can optimize for judge satisfaction without improving expert-grounded alignment, highlighting a Goodhart-like risk in judge--corrector feedback loops.
    \item We provide evidence that literature-oriented validation can enrich the disagreement space: calibrated with human validation and under strong assumptions (relationship isolation, a single human validator, and extrapolation), Deep Research yields substantial but bounded F1 gains; future work should test enrichment performance as these assumptions are relaxed. We also show a proof-of-concept judge-triggered deployment policy that improves cost-efficiency by selectively invoking Deep Research on uncertain edges.
\end{itemize}

%------------------------------------------------------------------------------
\subsection{Concluding Remarks}
\label{sec:conclusion_remarks}
%------------------------------------------------------------------------------

Taken together, these findings position LLMs as scalable assistants for proposing and stress-testing causal structures, not as autonomous validators. The most promising workflow is hybrid: use LLMs to expand the edge space, use judges and UQ signals to triage uncertainty, and selectively apply evidence-oriented validation (e.g., Deep Research with an explicit applicability threshold) to focus expert attention on the contested edges that most affect the final model.



The most significant finding may be the reconceptualization of ``hallucination'' itself in this context. The substantial proportion of LLM-generated edges classified as false positives against expert ground truth that nonetheless have retrievable causal evidence (62.4\%) suggests that the expert-LLM disagreement space contains overlooked-but-valid relationships alongside genuine fabrications. This reframing positions LLMs not merely as reproducers of expert consensus, but as tools for surfacing literature-supported causal relationships that experts may have missed---a function that becomes increasingly valuable as system complexity grows and the edge space ($O(|V|^2)$) exceeds the capacity for exhaustive expert deliberation.

As LLM capabilities continue to advance, the methods evaluated here---particularly the combination of uncertainty-guided selective deployment with literature-grounded validation---provide a foundation for integrating LLM-generated scientific knowledge into rigorous research workflows. The path forward requires not replacing expert judgment, but augmenting it: using LLMs to handle the scalable synthesis task while reserving expert attention for the specialized contributions that only domain knowledge can provide.




